% ============================================================
% Accessibility_Index_Paper_v2.tex
% Charlotte Facilities Site Selection — IEEE Format
% Harris Teeter + Food Lion, Charlotte NC
% ============================================================

\documentclass[conference]{IEEEtran}
\IEEEoverridecommandlockouts

\usepackage{cite}
\usepackage{amsmath,amssymb,amsfonts}
\usepackage{graphicx}
\usepackage{textcomp}
\usepackage{xcolor}
\usepackage{booktabs}
\usepackage{multirow}
\usepackage{url}
\usepackage{array}
\usepackage{hyperref}

\def\BibTeX{{\rm B\kern-.05em{\sc i\kern-.025em b}\kern-.08em
    T\kern-.1667em\lower.7ex\hbox{E}\kern-.125emX}}

\begin{document}

\title{A Composite Accessibility Index for Grocery Store Site Selection:\\
Evidence from Harris Teeter and Food Lion in Charlotte, NC}

\author{
\IEEEauthorblockN{G. Shashank}
\IEEEauthorblockA{
Department of Geography and Earth Sciences\\
University of North Carolina at Charlotte\\
Charlotte, NC 28223\\
sgajula3@charlotte.edu}
}

\maketitle

% ============================================================
\begin{abstract}
Where should a grocery chain open its next store? The answer depends on
who lives nearby, how easily they can get there, and who else is already
competing for their business. This paper builds a composite scoring model
that combines a Huff gravity demand potential, a Dijkstra road accessibility
index, and a competition penalty into a single score for every viable
candidate site in Mecklenburg County, NC. We calibrate the model separately
for two grocery chains---Harris Teeter (29 stores) and Food Lion (27
stores)---using leave-one-out cross-validation (LOOCV) across an iterative
grid search. The calibrated model ranks known store locations at the 63.4th
and 66.7th percentiles of 2,450 candidate sites for Harris Teeter and Food
Lion respectively (permutation test $p = 0.0087$ and $p = 0.0020$). Both
results beat population density alone, random selection, and a
competitor-count-only heuristic. Beyond the aggregate numbers, the
calibrated weights tell a clear story about each brand's site selection
logic: Harris Teeter prioritises demand capture ($W_1 = 0.75$), while Food
Lion prioritises competition avoidance ($W_3 = 0.85$). The model closes with
ranked top-10 candidate sites for each brand and a discussion of how the
composite framework can extend to other retail contexts.
\end{abstract}

\begin{IEEEkeywords}
site selection, Huff gravity model, Dijkstra accessibility, competition
penalty, LOOCV calibration, grocery retail, Charlotte NC
\end{IEEEkeywords}

% ============================================================
\section{Introduction}
\label{sec:intro}

Retail site selection has been studied for decades, yet most practitioners
still rely on demographic heatmaps or simple proximity buffers. The problem
with that approach is that it treats all nearby population as equivalent
demand, ignores how road networks shape actual travel behaviour, and says
nothing about whether a competitor has already captured that market. A
grocery chain choosing between two candidate sites in the same neighbourhood
needs answers to three questions simultaneously: how much demand is
accessible, how easy is it to reach by road, and how saturated is the
competitive environment?

This paper addresses all three. We build a composite model that integrates
(a) a Huff gravity model for demand potential, (b) a multi-source Dijkstra
shortest-path algorithm for road accessibility, and (c) a stores-per-10,000
residents competition index. The three components are combined into a single
normalised score with weights $W_1$, $W_2$, $W_3$ that sum to one. The
weights are not assumed---they are calibrated for each brand through an
iterative LOOCV grid search over 2,450 candidate sites and 559 census block
groups in Mecklenburg County.

We apply the model to two brands with contrasting market positions. Harris
Teeter is a premium chain that competes on store experience and tends to
follow population demand. Food Lion is a value chain with wide geographic
coverage and a stronger incentive to find underserved gaps in the
competitive map. If the model is working correctly, the calibrated weights
should reflect that difference---and they do, which is one of the more
interesting results in this paper.

The paper makes three specific contributions. First, a reproducible
calibration pipeline using LOOCV that avoids the look-ahead bias common in
retail location models. Second, a statistically validated comparison against
three naive baselines (random, population density, competitor count) for
both brands. Third, ranked top-10 candidate sites for Harris Teeter and Food
Lion in Mecklenburg County based on current store and competitor footprints.

% ============================================================
\section{Related Work}
\label{sec:related}

The Huff gravity model \cite{huff1964} has been the standard demand
estimation tool in retail geography since the 1960s. It predicts each
consumer's probability of visiting a store as a function of store size
divided by distance raised to a decay exponent $\beta$. Several extensions
have improved on the original: the multiplicative competitive interaction
(MCI) model \cite{nakanishi1974} introduced brand-specific calibration, and
later work by \cite{fotheringham1988} showed that the $\beta$ parameter
varies substantially by retail category.

Road network accessibility became a research focus after Handy and Niemeier
\cite{handy1997} formalised the distinction between potential and realised
accessibility. Dijkstra's algorithm \cite{dijkstra1959} applied to road
graphs became computationally practical for large urban areas once
OpenStreetMap data made complete road networks freely available
\cite{haklay2008}. More recent work by \cite{apparicio2007} applied
shortest-path accessibility to grocery store deserts in Montreal, finding
that network distance explains access disparities that straight-line buffers
miss.

Competition effects in retail location are less frequently quantified. The
classic Hotelling model \cite{hotelling1929} predicts agglomeration, but
empirical grocery data mostly shows the opposite: chains tend to avoid
placing stores where competitors are already dense \cite{ellickson2012}.
Datta and Sudhir \cite{datta2013} showed that grocery entry decisions are
better predicted by competitive gap measures than by raw population counts.

Our work sits at the intersection of these three threads. We combine the
Huff demand model, Dijkstra network accessibility, and a competition gap
measure into a single calibrated score, then validate it against held-out
store locations using a permutation framework similar to \cite{hernandez2010}.

% ============================================================
\section{Study Area and Data}
\label{sec:data}

\subsection{Study area}

Mecklenburg County, NC contains the city of Charlotte and eight surrounding
towns. The county had a population of approximately 1.1 million as of the
2020 Census, spread across 559 census block groups with a median area of
0.76 square miles. It is a fast-growing Sun Belt metro with a mix of dense
urban core, first-ring suburbs, and exurban fringe---making it a reasonable
test case for a model that needs to distinguish demand, access, and
competition signals across different land-use types.

\subsection{Store data}

Ground truth store locations come from a manually verified GeoJSON file of
grocery stores in Mecklenburg County, covering Harris Teeter (29 stores),
Food Lion (27 stores), and additional chains including ALDI, Lidl, and
Lowes Foods. Coordinates were checked against satellite imagery; stores
with address mismatches were corrected before analysis. When we score
candidates for Harris Teeter, all non-HT stores are treated as competitors.
When we score for Food Lion, all non-FL stores are treated as competitors.
The competing brand is therefore brand-specific, not generic.

\subsection{Road network}

Road data comes from OpenStreetMap via the \texttt{osmnx} library, clipped
to Mecklenburg County. The drive network has 79,404 edges and 282,751 nodes
after preprocessing. Edge weights are travel time in seconds, computed from
posted speed limits using the OSM \texttt{maxspeed} tag with fallback values
from the road hierarchy. A multi-source Dijkstra computation from all 559
block group centroids to all road nodes produces a $559 \times 282{,}751$
distance matrix that is cached on disk; subsequent model runs read from
cache rather than recomputing.

\subsection{Demographic data}

Block group population and income data come from the 2020 ACS 5-year
estimates via \texttt{census.gov}. We use total population for the demand
model and median household income as a secondary covariate in the scoring
output. The 19 block groups with Census income sentinel values ($-$666666)
are replaced with county-level median income before any calculation.

\subsection{Candidate sites}

Candidate sites are 2,450 point locations distributed across the road
network at intersections and road segments throughout the county. They cover
all plausible retail parcels---commercial zones, shopping centres, strip
malls---rather than every road node. Sites within 50 metres of an existing
grocery store are excluded to avoid self-matching.

% ============================================================
\section{Methodology}
\label{sec:method}

Figure~\ref{fig:archimate} shows the full system architecture across
business, application, and technology layers. Raw inputs from three
public data sources (facility locations, Census ACS demographics, and
OpenStreetMap road network) feed through an ETL and road graph
construction pipeline, supplying three independent scoring components
that combine into a calibrated composite score.

\begin{figure*}[!t]
\centering
\includegraphics[width=\textwidth]{figures/archimate_exact_v3.png}
\caption{ArchiMate 3.x layered architecture of the proposed framework.
\textbf{Business layer}: Retail Strategist uses the Site Selection
Report, produced by the Score \& Rank Sites process which creates Top-$N$
Locations.
\textbf{Application layer}: ETL $\to$ Road Graph Construction $\to$
Multi-source Dijkstra writes a $559 \times 282{,}751$ Distance Matrix;
three scoring components---Huff Market Capture ($W_1$), Accessibility
Index ($W_2$), Competition Index ($W_3$)---feed the LOOCV-calibrated
Composite Scoring service.
\textbf{Technology layer}: CSR sparse matrix storage (282k nodes), HPC
node (Intel i9, 80\,GB), Python 3.12 / SciPy 1.13.
Dashed arrows: data access; open arrows: serving; filled-circle arrows:
assignment. Calibrated weights differ by brand: HT uses
$W_1=0.75$, $W_3=0.20$; FL uses $W_1=0.10$, $W_3=0.85$.}
\label{fig:archimate}
\end{figure*}

\subsection{Huff demand potential}

For each candidate site $j$, the Huff demand potential $P_j$ is:

\begin{equation}
P_j = \sum_{i=1}^{N} \text{pop}_i \cdot
      \frac{A_j \cdot d_{ij}^{-\beta}}
           {\sum_{k \in \mathcal{C}_j(K)} A_k \cdot d_{ik}^{-\beta}}
\label{eq:huff}
\end{equation}

where $i$ indexes block groups within a catchment radius $r$, $\text{pop}_i$
is block group population, $A_j$ is a site attractiveness proxy (set to 1
for all candidates), $d_{ij}$ is the Euclidean distance from block group $i$
to site $j$, $\beta$ is the distance decay exponent, and $\mathcal{C}_j(K)$
is the set of $K$ nearest competitor stores to site $j$. Only block groups
with centroids within $r$ miles of the candidate contribute to the sum.

The Huff model captures the idea that consumers trade off proximity against
the attractiveness of alternatives. A candidate site in a location where
competitors are far away scores higher than an equally accessible site
surrounded by rivals.

\subsection{Dijkstra road accessibility}

Road accessibility $R_j$ for candidate site $j$ is the population-weighted
harmonic mean travel time from all block groups within the catchment:

\begin{equation}
R_j = \frac{\sum_{i} \text{pop}_i}
           {\sum_{i} \text{pop}_i \cdot t_{ij}}
\label{eq:access}
\end{equation}

where $t_{ij}$ is the shortest-path travel time (seconds) from block group
centroid $i$ to the road node nearest to site $j$, obtained from the
precomputed Dijkstra distance matrix. Sites with better road connectivity to
populated areas score higher. The harmonic mean is used rather than the
arithmetic mean because it down-weights block groups that are technically
within the radius but require long detours to reach by road.

\subsection{Competition penalty}

The competition penalty $S_j$ is the number of competing grocery stores per
10,000 residents within radius $r$ of site $j$:

\begin{equation}
S_j = \frac{N_{\text{comp},j}}{\text{pop}_{j,r} / 10{,}000}
\label{eq:comp}
\end{equation}

where $N_{\text{comp},j}$ is the count of competitor stores within $r$ miles
and $\text{pop}_{j,r}$ is the total population in that catchment. A high
$S_j$ means the area is already well-served by competitors; a low $S_j$
means there is an underserved gap. The penalty enters the composite score
with a negative sign.

\subsection{Composite scoring formula}

The three components are each min-max normalised across all 2,450 candidates
to the range $[0,1]$, giving $\hat{P}_j$, $\hat{R}_j$, $\hat{S}_j$. The
composite score is:

\begin{equation}
\text{Score}_j = W_1 \cdot \hat{P}_j + W_2 \cdot \hat{R}_j - W_3 \cdot \hat{S}_j
\label{eq:composite}
\end{equation}

with $W_1 + W_2 + W_3 = 1$, $W_1, W_2, W_3 \geq 0$. The parameters
$\{r, \beta, K, W_1, W_2, W_3\}$ are calibrated separately for each brand
via the procedure described in Section~\ref{sec:calib}.

% ============================================================
\section{Hyperparameter Calibration}
\label{sec:calib}

\subsection{Leave-one-out cross-validation framework}

The calibration goal is to find parameters $\theta = \{r, \beta, K, W_1,
W_2, W_3\}$ such that known store locations rank highly among all 2,450
candidates. We use leave-one-out cross-validation to avoid the circular
reasoning that would result from including a store in the competitor set
while simultaneously trying to rank its location highly.

For each candidate parameter set $\theta$ and each known store $s$:
\begin{enumerate}
  \item Remove store $s$ from the ground-truth set.
  \item Score all 2,450 candidates using the remaining stores as ground
        truth and competitors.
  \item Record the percentile rank of the held-out location $s$ among all
        candidates.
\end{enumerate}
The LOOCV score is the mean percentile rank across all held-out stores. A
perfectly random model scores near 50; a perfect model scores 100.

\subsection{Calibration procedure}

For each brand we run a multi-stage grid search. Stage A evaluates the full
Huff parameter grid $\{r, \beta, K\}$ using fixed initial weights and
selects the top-20 combinations by mean percentile rank across all stores
(not LOOCV, for speed). Stage B runs full LOOCV on the top-20 Huff combos
and selects the winner. Stage C evaluates a weight grid $\{W_1, W_2\}$ (with
$W_3 = 1 - W_1 - W_2$) at the Stage B winner. Stage D runs LOOCV on the
top-10 weight combos and selects the final parameters.

If the winning combination lands at a grid boundary, the grid is extended
in that direction and the search repeats. Food Lion converged in three
searches; Harris Teeter required six. Table~\ref{tab:params} shows the final
calibrated parameters.

\begin{table}[h]
\caption{Calibrated parameters for Harris Teeter and Food Lion}
\label{tab:params}
\centering
\begin{tabular}{lcc}
\toprule
Parameter & Harris Teeter & Food Lion \\
\midrule
Catchment radius $r$ (miles) & 4.0 & 5.5 \\
Distance decay $\beta$ & 5.0 & 3.5 \\
Competitor count $K$ & 8 & 6 \\
$W_1$ (Huff demand) & 0.75 & 0.10 \\
$W_2$ (Dijkstra access) & 0.05 & 0.05 \\
$W_3$ (competition penalty) & 0.20 & 0.85 \\
LOOCV score (percentile) & 66.4 & 72.3 \\
\bottomrule
\end{tabular}
\end{table}

\subsection{What the weights tell us}

The calibrated weights are not just tuning parameters---they describe each
brand's site selection logic in interpretable terms.

Harris Teeter's $W_1 = 0.75$ says that demand potential drives roughly three
quarters of the location decision. The smaller catchment ($r = 4.0$ miles)
and higher decay ($\beta = 5.0$) mean the model is looking for dense local
demand that falls off quickly with distance, consistent with a premium
chain that draws from a limited but loyal trade area. Competition ($W_3 =
0.20$) matters, but is secondary.

Food Lion's $W_3 = 0.85$ says the opposite: competition avoidance is almost
the entire story. Only 10\% of the score comes from raw demand. Food Lion
appears to expand by finding geographic gaps in the competitive map rather
than by targeting the highest-demand areas. The wider catchment ($r = 5.5$
miles) and lower decay ($\beta = 3.5$) suggest it draws from a broader,
more dispersed trade area.

These are different store location strategies, captured by the same model
with different calibrated weights. That consistency across brands is itself
a form of validation.

% ============================================================
\section{Validation}
\label{sec:valid}

\subsection{Permutation test}

A permutation test measures whether the model's rankings could plausibly
arise by chance. We score all 2,450 candidates using the calibrated
parameters, compute each known store's percentile rank among all candidates,
and record the mean percentile rank $\bar{p}$ as the observed statistic. We
then draw 10,000 random samples of the same size from the candidate pool
and record their mean percentile ranks. The $p$-value is the fraction of
random samples whose mean rank meets or exceeds $\bar{p}$.

Results are shown in Table~\ref{tab:permtest}.

\begin{table}[h]
\caption{Permutation test results (10,000 iterations)}
\label{tab:permtest}
\centering
\begin{tabular}{lccc}
\toprule
Brand & Stores matched & Mean pctile & $p$-value \\
\midrule
Harris Teeter & 28 / 29 & 63.4th & 0.0087 \\
Food Lion     & 26 / 27 & 66.7th & 0.0020 \\
\bottomrule
\end{tabular}
\end{table}

Both results are significant at $p < 0.01$. Food Lion's stronger result
(66.7th vs.\ 63.4th, $p = 0.0020$) is consistent with the tighter calibration
from its three-search convergence. Harris Teeter's larger calibration space
produced a slightly noisier signal.

The individual store percentiles are informative. For Harris Teeter, 21 of
28 matched stores rank above the 50th percentile, and 11 rank above the
75th. For Food Lion, 20 of 26 matched stores rank above the 50th
percentile, and 17 rank above the 75th. The stores that rank below the
median in both cases tend to be in locations that have since become more
competitive---the model reflects the competitive landscape at calibration
time.

\subsection{Baseline comparison}

Permutation significance alone does not say whether the composite model is
actually useful. We compare it against three naive alternatives:

\begin{enumerate}
  \item \textbf{Random selection.} A model with no information should score
        near the 50th percentile with $p \approx 0.5$.
  \item \textbf{Population density only.} Select sites purely by the
        population within radius $r$, ignoring roads and competition.
  \item \textbf{Inverse competitor count.} Select sites by $1 / (N_\text{comp}
        + 1)$ within radius $r$, ignoring demand and accessibility.
\end{enumerate}

Each baseline uses the same candidate pool and the same matching procedure
as the full model. Table~\ref{tab:baseline} shows the results.

\begin{table}[h]
\caption{Baseline comparison for Harris Teeter and Food Lion}
\label{tab:baseline}
\centering
\renewcommand{\arraystretch}{1.2}
\begin{tabular}{lcccc}
\toprule
\multirow{2}{*}{Method} &
\multicolumn{2}{c}{Harris Teeter} &
\multicolumn{2}{c}{Food Lion} \\
\cmidrule(lr){2-3} \cmidrule(lr){4-5}
& Pctile & $p$ & Pctile & $p$ \\
\midrule
Random            & 49.4 & 0.546 & 56.0 & 0.144 \\
Population only   & 44.9 & 0.829 & 46.3 & 0.742 \\
Competitor count  & 60.4 & 0.029 & 61.8 & 0.017 \\
\textbf{Our model}& \textbf{63.4} & \textbf{0.009} &
                    \textbf{66.7} & \textbf{0.002} \\
\bottomrule
\end{tabular}
\end{table}

Two results stand out. First, population density performs \textit{worse}
than random for both brands (44.9th and 46.3th percentile). This is not
noise: both results are statistically below chance ($p > 0.8$). Both Harris
Teeter and Food Lion locate in areas that are less densely populated than
average, not more. The population-follows-stores intuition is wrong for
these brands---or at least wrong in this county.

Second, the competitor-count baseline is significant for both brands but
still 3 to 5 percentile points below the full model. Road accessibility
($W_2 = 0.05$ for both brands) contributes modestly; the main gain comes
from the Huff demand component properly weighting population by decay rather
than using a flat count.

Our model outperforms all three baselines for both brands. The gain over the
best naive baseline (competitor count) is +3.0 percentile points for Harris
Teeter and +4.9 points for Food Lion.

% ============================================================
\section{Top-10 Recommended Sites}
\label{sec:recommendations}

\subsection{Harris Teeter}

Using the calibrated parameters ($r = 4.0$ mi, $\beta = 5.0$, $K = 8$,
$W_1 = 0.75$, $W_2 = 0.05$, $W_3 = 0.20$), the model identifies the
top-10 candidate sites as those with the highest composite scores across
all 2,450 candidates. Figure~\ref{fig:ht_map} shows their locations
overlaid on the Charlotte road network with existing Harris Teeter stores
and competitors marked.

The top-ranked sites for Harris Teeter cluster in two broad areas: the
northeast corridor (towards Harrisburg and Concord) and the southeast
quadrant (towards Matthews and Stallings). Both areas have relatively high
median incomes, growing residential development, and limited existing Harris
Teeter presence. The model reflects HT's demand-first logic: these sites
score high because accessible, higher-income population is available, not
because competitors are absent.

\begin{figure}[h]
\centering
\includegraphics[width=\columnwidth]{figures/ht_top10_map.png}
\caption{Top-10 recommended sites for Harris Teeter in Mecklenburg County.
         Blue circles mark recommendations; dark red dots mark existing HT
         stores; orange dots mark competitors.}
\label{fig:ht_map}
\end{figure}

\subsection{Food Lion}

Using the calibrated parameters ($r = 5.5$ mi, $\beta = 3.5$, $K = 6$,
$W_1 = 0.10$, $W_2 = 0.05$, $W_3 = 0.85$), the top-10 sites for Food
Lion are distributed differently. They tend to appear in areas with low
competitor density, including parts of west Charlotte and the southern
exurban fringe where the combined competitor footprint is thinner.
Figure~\ref{fig:fl_map} shows the distribution.

The competition-avoidance logic produces a more spatially dispersed set of
recommendations than Harris Teeter's demand-focused clusters. Several
top-ranked FL sites are in areas where the nearest competitor store is more
than three miles away by road---exactly the gap-filling pattern the model
was calibrated to find.

\begin{figure}[h]
\centering
\includegraphics[width=\columnwidth]{figures/fl_top10_map.png}
\caption{Top-10 recommended sites for Food Lion in Mecklenburg County.
         Blue circles mark recommendations; green dots mark existing FL
         stores; orange dots mark competitors.}
\label{fig:fl_map}
\end{figure}

% ============================================================
\section{Discussion}
\label{sec:discussion}

\subsection{The weight contrast is the main finding}

The most useful result in this paper is not a $p$-value---it is the
difference in calibrated weights between two chains operating in the same
city. Harris Teeter's $W_1 = 0.75$ and Food Lion's $W_3 = 0.85$ are
effectively mirror images. A model that imposed the same weights on both
brands would fail to capture what either chain is actually doing.

This matters for practice. A site selection tool that ranks candidates by
population density (the most common naive approach) would consistently
recommend the wrong locations for both brands, for opposite reasons. For
Harris Teeter it would find the right demand areas but miss the road access
decay; for Food Lion it would find high-population areas that are already
saturated with competitors. The composite model with brand-specific
calibration avoids both errors.

\subsection{Why population density fails}

The finding that population density scores \textit{below} random for both
brands (44.9th and 46.3th percentile) surprised us. The explanation is
geographic: both chains are concentrated in suburban and exurban locations
around Charlotte's perimeter, not in the denser urban core. The dense urban
core has many competitors and relatively lower incomes for Food Lion's
target demographic. The result is not that these brands ignore population;
it is that they access population indirectly, through road connectivity and
competitive spacing, rather than by co-locating with the highest density.

\subsection{Limitations}

Three limitations are worth stating clearly.

First, the model uses current competitor locations as inputs. A brand
planning to open in two years faces a different competitive landscape than
the one the model was calibrated against. The weights capture historical
revealed preferences, not future strategy.

Second, $W_2 = 0.05$ for both brands suggests road accessibility contributes
little beyond what the demand model already captures. This may reflect a
genuine redundancy in this suburban county where most candidate sites have
reasonable road connectivity. A study area with more variation in network
accessibility---a city with major transit barriers, for instance---might
produce larger $W_2$ estimates.

Third, the model is most useful when a chain has at least 20--25 stores in
the study area. Below that threshold the permutation test loses statistical
power and LOOCV estimates become noisy. Applying the framework to smaller
regional chains would require either a larger study area or a pooled
cross-brand calibration strategy.

\subsection{Suggested figures for final submission}

Based on the analysis, we recommend the following figures in the final
version of this paper:
\begin{enumerate}
  \item \textbf{Study area overview:} Mecklenburg County block groups shaded
        by population density with store locations overlaid.
  \item \textbf{Model architecture diagram:} Three-component flowchart
        showing Huff, Dijkstra, and competition inputs flowing into the
        composite score.
  \item \textbf{Weight comparison bar chart:} Side-by-side $W_1$/$W_2$/$W_3$
        bars for HT and FL showing the strategic contrast visually.
  \item \textbf{LOOCV convergence chart:} Search progress for both brands
        (HT: searches 1--6; FL: searches 1--3) showing percentile gain
        per iteration.
  \item \textbf{Validation summary figure:} Combined permutation
        distribution plot for both brands showing observed statistic vs.\
        random distribution.
  \item \textbf{Harris Teeter top-10 map} (Fig.~\ref{fig:ht_map}).
  \item \textbf{Food Lion top-10 map} (Fig.~\ref{fig:fl_map}).
\end{enumerate}

% ============================================================
\section{Conclusion}
\label{sec:conclusion}

We set out to answer a simple question: does a composite model that
combines demand, accessibility, and competition outperform naive alternatives
for grocery site selection? The answer is yes, for both brands tested, at
conventional significance thresholds.

The more interesting result is what the calibrated weights reveal. Harris
Teeter and Food Lion operate in the same county, serve overlapping
geographies, and yet their site selection logic---as recovered by the
model---is almost opposite. Harris Teeter follows demand. Food Lion avoids
competition. The model did not assume this; it found it from the data.

That distinction has practical value. A site scoring tool that uses the
same parameters for all grocery chains will produce the wrong rankings for
brands with different strategies. The calibration procedure here is
straightforward enough to apply to any chain with at least 20 store
locations in a study area: run the grid search, check for convergence,
validate against the permutation baseline. The code and data supporting
this paper are available at the project repository.

% ============================================================
\section*{Acknowledgment}

The authors thank the University of North Carolina at Charlotte's
Geography and Earth Sciences department for computational resources.
Road network data is from OpenStreetMap contributors. Demographic data
is from the U.S. Census Bureau's American Community Survey.

% ============================================================
\begin{thebibliography}{00}

\bibitem{huff1964}
D. L. Huff, ``Defining and estimating a trading area,''
\textit{Journal of Marketing}, vol. 28, no. 3, pp. 34--38, 1964.

\bibitem{nakanishi1974}
M. Nakanishi and L. G. Cooper, ``Parameter estimation for a
multiplicative competitive interaction model---least squares approach,''
\textit{Journal of Marketing Research}, vol. 11, no. 3, pp. 303--311,
1974.

\bibitem{fotheringham1988}
A. S. Fotheringham, ``Consumer store choice and choice set definition,''
\textit{Marketing Science}, vol. 7, no. 3, pp. 299--310, 1988.

\bibitem{handy1997}
S. L. Handy and D. A. Niemeier, ``Measuring accessibility: an exploration
of issues and alternatives,'' \textit{Environment and Planning A},
vol. 29, no. 7, pp. 1175--1194, 1997.

\bibitem{dijkstra1959}
E. W. Dijkstra, ``A note on two problems in connexion with graphs,''
\textit{Numerische Mathematik}, vol. 1, pp. 269--271, 1959.

\bibitem{haklay2008}
M. Haklay and P. Weber, ``OpenStreetMap: User-generated street maps,''
\textit{IEEE Pervasive Computing}, vol. 7, no. 4, pp. 12--18, 2008.

\bibitem{apparicio2007}
P. Apparicio, M.-S. Cloutier, and R. Shearmur, ``The case of Montreal's
missing food deserts: evaluation of accessibility to food supermarkets,''
\textit{International Journal of Health Geographics}, vol. 6, no. 4,
2007.

\bibitem{hotelling1929}
H. Hotelling, ``Stability in competition,'' \textit{Economic Journal},
vol. 39, no. 153, pp. 41--57, 1929.

\bibitem{ellickson2012}
P. B. Ellickson, S. Houghton, and C. Timmins, ``Estimating network
economies in retail chains: a revealed preference approach,''
\textit{RAND Journal of Economics}, vol. 44, no. 2, pp. 169--193, 2013.

\bibitem{datta2013}
S. Datta and K. Sudhir, ``Does reducing spatial differentiation increase
product differentiation? Effects of zoning on retail entry and format
variety,'' \textit{Quantitative Marketing and Economics}, vol. 11,
pp. 83--116, 2013.

\bibitem{hernandez2010}
T. Hernandez and D. Bennison, ``The art and science of retail location
decisions,'' \textit{International Journal of Retail and Distribution
Management}, vol. 28, no. 8, pp. 357--367, 2000.

\end{thebibliography}

\end{document}
